\chapter{Was die gesamte Sitzung ergeben hat}
\section{Das verständliche Gesamtbild}
Am Anfang stand die Hoffnung, dass eine sehr kleine geometrische Grammatik bereits das ganze Universum festlegt. Die Untersuchung hat dieses Ziel nicht erreicht. Sie hat aber aus einer allgemeinen Hoffnung eine wesentlich präzisere Kette von Fragen und nachprüfbaren Verbindungen gemacht.

Das passende Bild ist inzwischen ein Bauplan mit mehreren Übersetzungsstufen. Der Compiler legt ein Alphabet und Beziehungen zwischen seinen Zeichen fest. Ein Zustandsraum macht daraus mögliche Konfigurationen. Eine Dynamik sagt, wie sich diese Konfigurationen ändern. Lokale Felder sagen, welche Eingriffe an welchen Orten möglich sind. Erst eine gemeinsame räumliche Grenze könnte daraus die beobachtete Physik machen. Die Sitzung hat an mehreren Übergängen exakte Übersetzungen geliefert. Sie hat zugleich nachgewiesen, dass manche scheinbar passenden Übersetzungen wesentliche Information verlieren.

\finding{Der wichtigste gemeinsame Fortschritt}{Wir können heute wesentlich genauer unterscheiden, wann zwei Beschreibungen wirklich denselben Prozess darstellen. Übereinstimmen müssen nicht nur Zahlen oder Energien, sondern auch geladene Operatoren, Phasen, Statistik, Zustand und Zeit. Einige dieser gemeinsamen Abbildungen sind jetzt exakt konstruiert; ihre Auswahl aus den ursprünglichen Prinzipien ist noch nicht geschlossen.}

Ein einzelner kleiner Eigenwert war dafür ein schlechter Kompass. Er konnte die gemeinsame Umschaltung einer ganzen Quelle beschreiben, eine endliche Bindungsresonanz oder eine Bewegung über einem zu hohen Referenzzustand. Erst die zugehörigen Antwortoperatoren machten sichtbar, was tatsächlich angeregt wird. Ebenso war eine elegante E8-Zerlegung allein nicht genug: Ein entferntes Komplement konnte die fehlende Fermionparität, eine Lokalitätsphase, eine Determinantenwirkung oder die Ladungsabhängigkeit der Zeit tragen.

\section{Welche Ebenen jetzt miteinander verbunden sind}
\begin{longtable}{@{}P{.22\textwidth}P{.34\textwidth}P{.36\textwidth}@{}}
\toprule Ebene & Belastbarer Gewinn & Verbleibende Brücke\\\midrule\endfirsthead
\toprule Ebene & Belastbarer Gewinn & Verbleibende Brücke\\\midrule\endhead
Compiler und endliche Operatoren & Integrale Wörterbücher, native Ereignisse, Phasen- und Clockbedingungen sind in konkreten Darstellungen geprüft. & Ein Randpostulat bestimmt noch nicht automatisch die gesamte räumliche Energie und ihre Präparation.\\
Quantenquellen und Bindungen & Rückwirkung, überlappende Pakete und mehrere vollständige kleine Grundräume sind kontrolliert. Die früher verdeckte Bindungsantwort wurde identifiziert. & Große Phase, echte lokale Anregungen und die Herkunft der angenommenen Kopplungen bleiben offen.\\
Freie QWZ-Quelle & Ein lokaler geladener CAR-Grenzwert mit derselben Quellzeit ist für die wachsende Geometrie bewiesen. & Ein Kanal pro Rand erzeugt nicht von selbst acht Spezies, E8 oder die native Wechselwirkung.\\
Zusätzlicher zehnkanaliger Rand & Glue, Ladungen, Fermionparität, E8-Quotient und mögliche geladene Spinorfelder sind exakt bestimmt. & Zusatzkanäle, Energiematrix, kritischer Zustand und lokale Wechselwirkungen sind nicht aus P1/P2 ausgewählt.\\
Paar- und Sektorstruktur & Gemeinsame Intertwiner und einige Zustandsantworten sind konstruiert; nicht jeder passende Sektor liegt über demselben Vakuum. & Die native Fockalgebra und die zusammengesetzten Randfelder brauchen dieselbe Gruppen-, Zeit- und Lokalitätswirkung.\\
Vierphasenregister & Eine korrelierte Kodierung erhält volle CAR, Registerwechsel und Zeit. Produkt-und-Trace sowie korrelierte Kodierung sind entschieden verschieden. & Die feste Randholonomie erzeugt im kontrollierten lokalen Bulk-Limes keine neue Kraft.\\
Raumzeit und Weltphysik & Präzise gemeinsame Abnahmekriterien und mehrere Ausschlüsse falscher Abkürzungen. & Dynamische Geometrie, 3+1D-Chiralität, Kopplungsauswahl, kosmologischer Zustand und universelle Gravitation sind nicht hergeleitet.\\\bottomrule
\end{longtable}

Diese Zeilen gehören zu unterschiedlichen, ausdrücklich benannten Konstruktionen. Der endliche Wurzelraum, die freie QWZ-Quelle, das native 64-CAR/60-CCR-Modell, eine affine E8-Algebra und der zusätzliche zehnkanalige Rand sind nicht schon wegen gemeinsamer Zahlen dieselbe Theorie. Genau ihre gemeinsamen Abbildungen wurden geprüft. Die Tabelle ist deshalb keine Sammlung von Eigenschaften, die man ohne Weiteres zu einer fertigen Welt zusammensetzen könnte.

\section{Was wir tatsächlich besser verstanden haben}
\textbf{Erstens: Die Quelle muss ihre eigene Rückwirkung behalten.} Eine unveränderliche Quelle kann eine schöne Ereignisstatistik erzeugen, aber das folgende Ereignis hängt davon ab, was nach dem ersten mit der Quelle geschieht. Die Rückwirkungsmodelle und Registerrechnungen zeigen konkret, an welcher Stelle diese Information verloren geht oder erhalten bleibt. \src{08}\src{17}\src{51}\src{52}

\textbf{Zweitens: Bindungsänderungen sind echte Freiheitsgrade.} Die Koordinaten einer Wurzel sprechen nicht jede niedrige Anregung an. Für das Quellenpaar existiert ein kontrollierter Bindungsmodus, den lineare Wurzelkoordinaten durch Symmetrie nicht sehen. Beim Vergrößern darf man diesen Modus jedoch nicht unverändert kopieren: Das Vakuum kann seine Bindungsordnung neu organisieren. Die Domänenwandidee bleibt in geeigneten Klassen mathematisch konkret; ihre nackte kritische Paketfortsetzung ist als vollständiges Vakuum ausgeschlossen. \src{24}\src{26}\src{28}\src{29}

\textbf{Drittens: Geladene Felder prüfen die gemeinsame Zeit strenger als neutrale Ströme.} Eine Viertelholonomie kann aus neutralen Antworten verschwinden und dennoch Paar- und Spinorfrequenzen verändern. Der wachsende lokale QWZ-Grenzwert ist ein positiver gemeinsamer Feld-Zeit-Anschluss unter klaren Annahmen. Er löst nicht nachträglich die verschiedenen Festkreis- und Mehrkanalprobleme. \src{32}\src{33}

\textbf{Viertens: Das Komplement ist Teil des physikalischen Wörterbuchs.} Im lokalen Randgitter kann es die Fermionparität und die zur Lokalität nötige Phase tragen. Im Quellfunktional kann es einen Determinantenfaktor liefern. Ein gemeinsamer Faktor verschwindet aber aus einer relativen Down-/Lepton-Phase wieder. Er muss seine unterscheidbare Wirkung aus derselben Quelle erhalten. \src{37}\src{39}

\textbf{Fünftens: Erhaltene Information ist noch keine ausgewählte Naturdynamik.} Die jüngste Registerkodierung ist eine exakte Verbindung von Feldern und Zeit. Dennoch bleibt sie in jedem festen Teilchenzahlsektor eine freie Bewegung mit Randphase. Eine bessere Darstellung derselben Quelle kann verdeckte Übergänge aufdecken, aber keine zusätzliche lokale Streuwechselwirkung herbeizaubern. \src{52}

\section{Warum die Gesamtlösung weiterhin fehlt}
Es fehlt nicht nur ein letzter passender Zahlenwert. Die erste offene Auswahl liegt schon davor: Welche vollständige physische Quelle und welche Dynamik folgen aus den ursprünglichen Postulaten? Die Sitzung besitzt mehrere gute Kandidaten und exakte Teilverbindungen, aber noch keinen gemeinsamen Satz, der ihre unabhängigen Entscheidungen beseitigt.

Ein vereinfachtes Beispiel zeigt die Bedeutung. Aus denselben zulässigen Zuständen lassen sich unterschiedliche positive Energien bilden. Sie können dieselbe Symmetrie, dieselben Ladungen und sogar denselben Grundzustand besitzen, aber verschiedene Anregungsgeschwindigkeiten oder Streuprozesse erzeugen. Ein noch so genauer Zustandstest entscheidet zwischen ihnen nicht. Dafür braucht man den Generator beziehungsweise ein gleichwertiges vollständiges Mehrzeitgesetz. \src{STATE}

Auch nach einer solchen Auswahl wären große Schritte nötig: Der relevante Zustand müsste auf größeren Systemen existieren; seine lokalen Antworten müssten den richtigen Grenzwert besitzen; die geladenen Fermionfelder dürften nicht durch den gewählten Quotienten verschwinden; und der räumliche Verbund müsste selbst begründet werden. Erst dann wären vierdimensionale Chiralität, Spiegelentkopplung und eine universelle gravitative Antwort am selben Modell zu prüfen.

Die Untersuchung beweist weder, dass eine vollständige TFPT-Ursprungstheorie unmöglich ist, noch dass ihre Lösung schon verborgen vorliegt. Sie entscheidet konkrete Strukturen. Ihre stärkste Grenze ist zugleich eine positive Arbeitsanweisung: Die nächste Konstruktion muss die noch getrennten Voraussetzungen auf derselben Quelle zusammenführen, statt Eigenschaften aus verschiedenen Modellen nebeneinanderzustellen.

\section{Die Rolle von Einfachheit und Eleganz}
Die einfachen Ideen der Sitzung waren besonders produktiv, wenn sie eine entscheidende Annahme sichtbar machten: die Differenz von Summe und relativer Bewegung, ein einziger Zentrumstest, der gemeinsame Determinantenfaktor oder ein Register, das mit der Teilchenzahl weiterdreht. Diese kleinen Tests ersparten große Rechnungen an falschen Identifikationen.

Eleganz ist dabei ein Suchprinzip, kein Beweis dafür, dass das gesuchte Ergebnis existieren muss. Die kleinste tragfähige Theorie müsste nicht möglichst wenig Mathematik enthalten, sondern möglichst wenige unabhängige physikalische Entscheidungen. Der erzielte Fortschritt lässt sich daran messen, welche dieser Entscheidungen wirklich abgeleitet, welche gemeinsam gebunden und welche nur präziser als offene Annahme erkennbar wurden.

\chapter{Die größten anschließenden Forschungsfragen}
\section{Priorität 1: die Quelle wählt ihre Dynamik}
\textbf{Frage.} Welche einzige, aus den Originalprinzipien begründete Quelle erzeugt die benötigte nichtquadratische Bewegung und die gemeinsame Energie der geladenen Felder?

Die freie QWZ-Quelle, der native Paar-Fock-Hamiltonoperator und der zehnkanalige Rand enthalten jeweils andere tatsächlich gewählte Struktur. Die Registeruntersuchung hat eine konkrete Hoffnung entschieden: Eine korrelierte Hebung bewahrt zwar mehr Information als ein Trace, erzeugt im kontrollierten räumlichen Innenbereich aber keine neue lokale Wechselwirkung. Die Energiematrizen des Randkandidaten dürfen ebenfalls nicht als bloße Basiswechsel behandelt werden.

\textbf{Nächster entscheidender Nachweis.} Einen nativen Operator oder ein vollständiges Quellfunktional angeben und zeigen, wie daraus die benötigten Kanäle, die Energie sowie die relevante nichtlineare Kopplung entstehen. Wo ein Zusatzsektor gebraucht wird, ist seine Herkunft Teil des Nachweises. Falls eine kompakte Eichverbindung die Brücke sein soll, sind ihr Maß und ihre tatsächliche Flussamplitude herzuleiten; der bereits bekannte Index allein genügt nicht.

\textbf{Erfolgskriterium.} Ein gemeinsames Operator-/Zustands-/Zeitdiagramm ohne separat eingesetzte relative Kinetiken und ohne frei gewählte Reparaturenergie. \textbf{Abbruchkriterium für einen Kandidaten.} Wenn dessen Kernwirkung unter der eigenen Reduktion frei bleibt oder eine benötigte Gruppenwirkung exakt verschwindet, wird dieser Kandidat ohne neue begründete Voraussetzung nicht weiter als Ursprung der fehlenden Wechselwirkung optimiert.

\section{Priorität 2: geladene Fermionen und neutrale Ströme in einem lokalen Raum}
\textbf{Frage.} Wie tragen dieselben lokalen Zustände E8-Ströme, ungerade geladene Felder, die tatsächlichen Clocks und ihre Zeitentwicklung?

Die Antwort kann nicht nur ein neutraler Quotient sein. Im nativen Fockmodell sperrt ein Zentrumselement die direkte Operatoridentifikation. Im Randmodell rekonstruiert eine nichtlineare Stromwirkung die Spinorfelder, benötigt aber eine zusätzlich ausgewählte Energie. Das gerade Nachbargitter repariert den vollen Stromlift und entfernt dafür ursprüngliche ungerade Felder aus der lokalen Algebra.

\textbf{Nächster Nachweis.} Einen konkreten geladenen Sektorintertwiner mit Ladung, Parität, Kokzyklus, Trägerbereich und Zeit angeben. Danach seine Produkte mit den neutralen Strömen und mindestens einen Mehrzeitkorrelator am selben Zustand vergleichen. Ein Wilsonlinien-Endpunkt oder Twistfeld ist zulässig, sofern seine Lokalitäts- und Sektorenregel ausdrücklich mitgeliefert wird.

\textbf{Erfolg.} Nichtverschwindende gemeinsame Wirkung auf den tatsächlich verlangten Feldern, einschließlich der zentralen Zeichen. \textbf{Abbruchkriterium.} Ein falsches Zentrumzeichen, eine nicht entfernbare Monodromie oder ein nichtlokaler Rest widerlegt die beanspruchte lokale Identifikation. Ein Gleichsetzen von Dimensionen oder Labels ersetzt diesen Test nicht.

\section{Priorität 3: eine relative Flavorphase aus demselben Funktional}
\textbf{Frage.} Warum erzeugt die ursprüngliche Quelle unterschiedliche Down- und Lepton-Deformationen mit einer physisch unterscheidbaren relativen Phase?

Hier treffen die eingebrachte Lane-Synthese und die eigene Randprüfung direkt zusammen. Der vollständige Determinantenfaktor muss erhalten bleiben. Ein gemeinsamer massiver Faktor kürzt sich aber aus der relativen Antwort heraus. Eine konstante freie Dirac-Massenphase ist unitär entfernbar. Der native antisymmetrische Familienblock ist zudem zunächst rangdefizient; eine abstrakte Rangergänzung ist noch kein erlaubter lokaler Yukawa-Term des gewählten niedrigen Sektors.

\textbf{Nächster Nachweis.} Beide Deformationen im gleichen, bereits ausgewählten Quellfunktional definieren und ihre relative Antwort einschließlich Komplement, Ladung und Parität berechnen. Der Unterschied muss aus einem Quelloperator folgen, nicht aus zwei unabhängig eingesetzten Parametern.

\textbf{Erfolg.} Eine basisunabhängige, kanalunterscheidende Phase und der benötigte Rang am gemeinsamen physischen Zustand. \textbf{Abbruchkriterium.} Hebt sich die Phase in der vollen Determinante auf oder ist die vorgeschlagene Mischung paritätsverboten, ist dieser Mechanismus entschieden; weitere Fits an der reduzierten Matrix ändern daran nichts.

\section{Priorität 4: eine große Phase mit wirklicher lokaler Bewegung}
\textbf{Frage.} Gibt es in der physisch begründeten Hamiltonfamilie einen stabilen großen Zustandszweig, über dessen tatsächlichem Vakuum lokale geladene Anregungen linear dispergieren?

Die kleine Quellenlücke, die exakt laufende freie Vergleichsanregung und die nackte Ising-Paketkette beantworten diese Frage nicht. Die letzte Familie ist auf ihrer ganzen behaupteten kritischen Linie energetisch ausgeschlossen. Eine dressierte gemischte Phase ist damit nicht allgemein widerlegt, braucht aber eine kontrollierte Rückwirkung der herausprojizierten Zustände.

\textbf{Nächster Nachweis.} Vakuumvergleich, räumliche Spektralfunktion und Komplementabschätzung gemeinsam durchführen. Die Reduktionsfehler müssen gegenüber der mit der Größe relevanten Energieskala kontrolliert bleiben; ein kleiner Fehler bei fester Kettenlänge genügt nicht.

\textbf{Erfolg.} Derselbe Zustand trägt die niedrigste Energie, eine Folge lokaler Wellenzahlen, nichtverschwindendes geladenes Antwortgewicht und einen kontrollierten $z=1$-Grenzwert. \textbf{Abbruchkriterium.} Ein tieferer Konkurrenzzustand, extensiver unkontrollierter Rest oder ein verschwindender lokaler Übergang schließen den jeweils behaupteten Zweig aus.

\section{Priorität 5: Instrument, Anfangszustand und räumliche Beziehungen}
\textbf{Frage.} Wie werden nicht nur Zustände und Energien, sondern die tatsächlichen Ereignisse, ihre Records und ihre Orte aus demselben Prozess bestimmt?

Ein Grundzustand ist keine autonome Präparation. Eine POVM-Wahrscheinlichkeit bestimmt keinen Nachzustand. Ein festes Graphlabel erzeugt keine Dynamik der Inzidenz. Die Sitzung enthält konkrete Zeugen für alle drei Unterschiede. Die blockdiagonale Erweiterung des Bindungsmodells bevorzugt ohne weitere Regel getrennte Matchings und liefert keinen zusammenhängenden Raum.

\textbf{Nächster Nachweis.} Zunächst ein aus der Quelle festgelegtes Mehrzeitinstrument und zwei operational unterscheidbare lokale Algebren angeben. Anschließend einen tatsächlichen Übergang zwischen Bindungs- oder Inzidenzmustern herleiten, sofern daraus die Raumstruktur entstehen soll. Die originale Auslesestatistik muss auf diesem Prozess wiedergewonnen werden.

\textbf{Erfolg.} Zustandsvorbereitung, lokale Eingriffe, Ereignisgedächtnis und begrenzte Ausbreitung passen zu demselben Generator. \textbf{Abbruchkriterium.} Wenn zwei angeblich verschiedene Orte durch alle maßgeblichen Observablen identifiziert werden oder Folgeeingriffe vom reduzierten Zustand nicht bestimmbar sind, fehlt die beanspruchte operationale Rekonstruktion.

\section{Priorität 6: gemeinsame 3+1-dimensionale Physik}
\textbf{Frage.} Trägt die erfolgreiche Quelle anschließend chirale Materie, die gemessenen Kopplungen und einen universell gekoppelten masselosen Spin-2-Sektor?

Interne Spinorindizes sind keine Raumzeitspinoren. Eine 1+1-dimensionale Anomaliebilanz ist keine vierdimensionale Standardmodell-Realisierung. Ein linearer Zweig beweist keine gemeinsame Lichtgeschwindigkeit. Diese Unterscheidungen sind Voraussetzungen für den nächsten Nachweis, keine nachträglichen Vorbehalte.

\textbf{Erfolg.} Ein gemeinsamer lokaler Parent mit kontrolliertem Kontinuum, Chiraliät und Spiegelbehandlung, passender Ladungsstruktur, relativen Kopplungen, universeller gravitativer Antwort sowie begründetem kosmologischem Zustand. Frühere TFPT-Zahlenbeziehungen müssen daraus als Auslesungen folgen. Ein Ergebnis aus jedem von sechs verschiedenen Hilfsmodellen erfüllt diesen Vertrag nicht.

Die Arbeitsreihenfolge beginnt deshalb mit der Herkunftsfrage aus Priorität 1 und prüft daran Prioritäten 2 und 3. Große Phasen, räumliche Rekonstruktion und vierdimensionale Physik werden auf demselben Anschluss entwickelt. Die Liste verlangt keine sechs neuen Modellzweige.

\section{Dokumentarische Nacharbeiten}
Zwei frühere Aussagen benötigen unabhängig davon eine Bereinigung: der vollständige Drei-Quellen-Komplementboden und die allgemeine Vierersektor-Leakageformel. In diesem Dossier bleiben die nicht wiederhergestellten Teile ausdrücklich außerhalb der gesicherten Vollsätze. Ihre Rekonstruktion würde den Bestand verbessern, ist aber von einem neuen Fortschritt zur Weltphysik zu unterscheiden.

\chapter{Provenienz, Reproduzierbarkeit und Reichweite}
\section{Abdeckung und Quellenstand}
Diese zweite Fassung erfasst die gesamte in der Sitzung dokumentierte Ergebnisfolge vom 18. bis 20. September 2026 bis einschließlich der geladenen Registerrekonstruktion. Grundlage sind 55 vorherige Abschlussantworten, die vom Nutzer eingebrachten Anschlussberichte und die zugehörigen lokalen Forschungsartefakte. Die bereits vorhandene erste Fassung deckte die ersten 43 Abschlussantworten ab. Sie bleibt unverändert erhalten.

Die Darstellung gruppiert wiederholte Fortsetzungen nach ihrem mathematischen Inhalt. Sie ist eine vollständige thematische Zusammenfassung, kein wörtliches Gesprächsprotokoll. Programmquelltexte, große Matrizen und sämtliche einzelnen Prüfzeilen sind über die Originalverträge auffindbar; ihre vollständige Wiedergabe würde die Gesamterklärung unlesbar machen. Neu ist hier die redaktionelle Zusammenführung, kein zusätzlicher Forschungsbeweis.

Der Repository-HEAD beim Abgleich lautet
\begin{center}\small\ttfamily 66b91e40e245569f06ab440ead80f446c9be0ee5.\end{center}
Die Forschungsergebnisse liegen zusätzlich im Arbeitsbaum. Der Commit allein identifiziert deshalb nicht den dokumentierten Stand. Die begleitende Datei mit der Endung \texttt{.sources.json} bindet die verwendeten Dateien mit SHA-256, ihrer Größe und ihrer Rolle. Alle 95 Quellenpins der ersten Fassung waren beim aktuellen Abgleich vorhanden und bytegleich. Zusätzliche Quellen der Fortsetzungen sind gesondert erfasst.

S01--S29 bezeichnen die ursprüngliche nummerierte Ergebnisfolge. Sie entspricht 28 Vertragsindizes, weil .04 und .05 einen Vertrag teilen und .10 anders benannt ist. S30--S52 ordnen die späteren Quellen-, Rand-, Paar- und Registerarbeiten zu. Diese redaktionellen S-Nummern sind keine neue theoretische Claimnummerierung und keine neue Freigabe des Korpus.

Der Theoriegraph hilft, Verträge und Abhängigkeiten zu finden. Er ist keine Beweisquelle. Ebenso gelten eingereichte Lane-Synthesen als Hinweise auf prüfbare Anschlüsse; die gemeinsame physische Realisierung folgt nicht aus ihrer Verlinkung.

\section{Was geprüft wurde und was nicht}
Für diese zweite Dokumentation wurden Quellenstände, Vertragszuordnungen, Korrekturen und die Abdeckung der späteren Sitzung abgeglichen. LaTeX-Kompilation, interne Verweise und die gerenderte Seitendarstellung werden gesondert geprüft. Die großen numerischen und exakten Forschungsläufe wurden zur PDF-Erstellung nicht vollständig wiederholt. Ihre Zahlen bleiben als dokumentierte Resultate der jeweils bezeichneten ursprünglichen Rechnung ausgewiesen.

Die erste Fassung hatte zusätzlich einen Repository-Audit mit Exitcode null dokumentiert, einschließlich 266 bestandener RH-Prüfungen. Das war ein historischer Integritätslauf am 20.09.2026, kein neuer Lauf dieser zweiten Fassung und kein Beweis der physikalischen Schließung. Der damalige schnelle Modus übersprang Module und meldete erklärte Lean-\texttt{sorry}-, Graphzyklus- und Quellen-Drift-Warnungen. Das zugehörige Protokoll bleibt im Quellenregister gebunden.

Bei widersprüchlichen Formulierungen wird der präzisere spätere Belegstand verwendet. Deshalb stehen etwa der geschlossene Vierquellenvergleich und der nur teilweise wiederhergestellte Dreiquellensatz mit unterschiedlichen Status im selben Dokument. Eine unabhängige Begutachtung oder Formalisierung aller Forschungsbeweise wird nicht behauptet.

Die LaTeX-Datei ist eigenständig. In einer passenden TeX-Live-Umgebung kann sie mit \texttt{latexmk -pdf} erneut gesetzt werden. Das Quellenmanifest dokumentiert die finale Ausgabe und ihre Prüfungen. Ein sauber gesetztes PDF erhöht die Zugänglichkeit und Nachvollziehbarkeit; es ersetzt keine theoretische Herleitung.

\section{Notation und typische Verwechslungen}
\begin{longtable}{@{}P{.22\textwidth}P{.7\textwidth}@{}}
\toprule Zeichen & Bedeutung im jeweiligen Abschnitt\\\midrule\endfirsthead
\toprule Zeichen & Bedeutung im jeweiligen Abschnitt\\\midrule\endhead
$J$ in $H$ & Nichtnegative Ereigniskopplung des Quellenkettenmodells.\\
$J$ als Matrix & Viertelclock; weder die Kopplung noch eine Kodierungsabbildung.\\
$J_s,J_s^u$ & Isometrien der geladenen Registerkodierung.\\
$J_n^a$ & Affiner Strommodus.\\
$C$ & Tatsächliche Coxeterwirkung; $C_R$ in der Registerrechnung ist eine andere, kombinierte Erhaltungsgröße.\\
$\mathcal C_{60}$ & Kanal aus 60 Quellenereignissen, kein Coxeterelement.\\
$K$ & Im Randkapitel die Gitterform $I_{9,1}$; bei Quellenphasen gegebenenfalls ein Gesamtphasensektor.\\
$g$ & Je nach expliziter Definition Gruppenoperation, Kopplungszahl oder Ladungsvektor.\\
$q$ & Im Gitterteil globale Markierung; $q_j\in\{0,1\}$ im Registerteil markiert nur den Phasenbogen.\\
$\Delta$ & Energielücke bei einem benannten Hamiltonoperator oder Skalendimension eines bezeichneten Feldes; diese Größen sind nicht austauschbar.\\
$m,N$ & Kettenlänge, Registerzahl oder Teilchenzahl nur nach der jeweiligen lokalen Definition.\\
$P,Q=I-P$ & Gewählter Teilraum und Komplement. $PHP$ allein ist ohne Invarianz oder Fehlerkontrolle keine Volltheorie.\\
$F,T$ & Im Randkapitel verschiedene Gitterschnitte; $F_j$ im Registerkapitel ein Fermionfeld.\\\bottomrule
\end{longtable}


\appendix


\chapter{Quellen- und Ergebnisregister}
Alle Kennungen S01--S52 verweisen auf die zugehörigen Originalverträge. Maßgeblich sind deren konkrete Teilverdicts und Voraussetzungen; das gemeinsame Gesamtverdict lautet jeweils \texttt{PARTIAL}. Ein Link auf eine Beweisdatei ist keine Behauptung unabhängiger Begutachtung. Die vollständigen SHA-256-Werte aller hier gebundenen Dateien stehen im begleitenden Manifest.

\textbf{Pfadwurzeln.} \texttt{TFPT/} bezeichnet das Repository \path{/Users/stefanhamann/Projekte/tfpt-theoryv4/}. \texttt{SESSION/} bezeichnet \path{/Users/stefanhamann/Documents/Codex/2026-09-18/unte-2/}. Diese Präfixe machen die Fundstellen lesbar; das JSON-Manifest enthält die absoluten Pfade.

\section{Die nummerierte Ergebnisfolge}
\subsection*{S01 — Integrale Trialität}\label{s:01}
{\small\path{UR.COMPILER.INTEGRAL_TRIALITY.01}\par}
Vertrag: \path{TFPT/experiments/theory-contracts/compiler-integral-triality-20260918/contract_index.json}\par{\footnotesize SHA-256 (Anfang): \texttt{54d6580e9413d4c8}.}\par
Herleitung: \path{TFPT/experiments/theory-contracts/compiler-integral-triality-20260918/PROOF.txt}\par{\footnotesize SHA-256 (Anfang): \texttt{123e076cff29adca}.}\par
\subsection*{S02 — Starrheit der integralen Trialität}\label{s:02}
{\small\path{UR.COMPILER.TRIALITY_RIGIDITY.02}\par}
Vertrag: \path{TFPT/experiments/theory-contracts/compiler-triality-rigidity-20260918/contract_index.json}\par{\footnotesize SHA-256 (Anfang): \texttt{43da16e1aa0762bb}.}\par
Herleitung: \path{TFPT/experiments/theory-contracts/compiler-triality-rigidity-20260918/PROOF.txt}\par{\footnotesize SHA-256 (Anfang): \texttt{8cbc553eb0427d42}.}\par
\subsection*{S03 — Natives Wurzel- und CAR-Wörterbuch}\label{s:03}
{\small\path{UR.COMPILER.NATIVE_ROOT_DICTIONARY.03}\par}
Vertrag: \path{TFPT/experiments/theory-contracts/compiler-native-root-dictionary-20260918/contract_index.json}\par{\footnotesize SHA-256 (Anfang): \texttt{38c3007af98fafff}.}\par
Herleitung: \path{TFPT/experiments/theory-contracts/compiler-native-root-dictionary-20260918/PROOF.txt}\par{\footnotesize SHA-256 (Anfang): \texttt{b5de7459d2c849ea}.}\par
\subsection*{S04 — E6-Auslese und Determinantenergänzung}\label{s:04}
{\small\path{UR.COMPILER.E6_READOUT_CLOSURE.04}\par}
Vertrag: \path{TFPT/experiments/theory-contracts/compiler-e6-seam-readout-20260918/contract_index.json}\par{\footnotesize SHA-256 (Anfang): \texttt{6478c8ffd27d7901}.}\par
Herleitung: \path{TFPT/experiments/theory-contracts/compiler-e6-seam-readout-20260918/PROOF.txt}\par{\footnotesize SHA-256 (Anfang): \texttt{a40b401a71eca7a0}.}\par
\phantomsection\label{s:05}\textbf{S05 — Nahtcharakter auf E8.} \path{UR.COMPILER.SEAM_E8_CHARACTER_EXTENSION.05} ist ein eigenständiges verwandtes Resultat desselben Vertrags und derselben Beweisdatei; kein fehlender fünfter Beweisordner.

\subsection*{S06 — Affine Vakuumantwort und E6-Kubik}\label{s:06}
{\small\path{UR.COMPILER.VACUUM_CURRENT_CUBIC.06}\par}
Vertrag: \path{TFPT/experiments/theory-contracts/compiler-vacuum-current-cubic-20260918/contract_index.json}\par{\footnotesize SHA-256 (Anfang): \texttt{32daea34ae13bd70}.}\par
Herleitung: \path{TFPT/experiments/theory-contracts/compiler-vacuum-current-cubic-20260918/PROOF.txt}\par{\footnotesize SHA-256 (Anfang): \texttt{d65189b806facbb1}.}\par
\subsection*{S07 — Quellenrang und Clockwörterbuch}\label{s:07}
{\small\path{UR.COMPILER.SOURCE_CHANNEL_GATE.07}\par}
Vertrag: \path{TFPT/experiments/theory-contracts/compiler-source-channel-gate-20260918/contract_index.json}\par{\footnotesize SHA-256 (Anfang): \texttt{235cc326022476b1}.}\par
Herleitung: \path{TFPT/experiments/theory-contracts/compiler-source-channel-gate-20260918/PROOF.txt}\par{\footnotesize SHA-256 (Anfang): \texttt{7c9e47e12d06c94c}.}\par
\subsection*{S08 — Operations- und Prozessrekonstruktion}\label{s:08}
{\small\path{UR.COMPILER.PROCESS_RECONSTRUCTION.08}\par}
Vertrag: \path{TFPT/experiments/theory-contracts/compiler-process-reconstruction-20260918/contract_index.json}\par{\footnotesize SHA-256 (Anfang): \texttt{baa1c6d47d1a955b}.}\par
Herleitung: \path{TFPT/experiments/theory-contracts/compiler-process-reconstruction-20260918/PROOF.txt}\par{\footnotesize SHA-256 (Anfang): \texttt{c624193a3fde4c2b}.}\par
\subsection*{S09 — Quartischer Übertrag}\label{s:09}
{\small\path{UR.COMPILER.QUARTIC_LIFT.09}\par}
Vertrag: \path{TFPT/experiments/theory-contracts/compiler-quartic-carry-20260919/contract_index.json}\par{\footnotesize SHA-256 (Anfang): \texttt{0916af9829945938}.}\par
Herleitung: \path{TFPT/experiments/theory-contracts/compiler-quartic-carry-20260919/PROOF.txt}\par{\footnotesize SHA-256 (Anfang): \texttt{e38c31499b54bab6}.}\par
\subsection*{S10 — Lokalisierter Spinorsektor}\label{s:10}
{\small\path{UR.SEAM.LOCALIZED_SECTOR.10}\par}
Vertrag: \path{TFPT/experiments/theory-contracts/source-localized-spinor-20260919/contract_index.json}\par{\footnotesize SHA-256 (Anfang): \texttt{1302407ab999b263}.}\par
Herleitung: \path{TFPT/experiments/theory-contracts/source-localized-spinor-20260919/PROOF.txt}\par{\footnotesize SHA-256 (Anfang): \texttt{2a08b4d68fc27030}.}\par
\subsection*{S11 — Lokaler Stromursprung und neutraler Zellentest}\label{s:11}
{\small\path{UR.COMPILER.LOCAL_CURRENT_ORIGIN.11}\par}
Vertrag: \path{TFPT/experiments/theory-contracts/compiler-local-current-origin-20260919/contract_index.json}\par{\footnotesize SHA-256 (Anfang): \texttt{9cad6bb44b631fc8}.}\par
Herleitung: \path{TFPT/experiments/theory-contracts/compiler-local-current-origin-20260919/PROOF.txt}\par{\footnotesize SHA-256 (Anfang): \texttt{5270d1f86ae6c3d1}.}\par
\subsection*{S12 — Quartischer Phasentransport}\label{s:12}
{\small\path{UR.COMPILER.QUARTIC_HOLONOMY.12}\par}
Vertrag: \path{TFPT/experiments/theory-contracts/compiler-quartic-holonomy-20260919/contract_index.json}\par{\footnotesize SHA-256 (Anfang): \texttt{3bd9449d2f437753}.}\par
Herleitung: \path{TFPT/experiments/theory-contracts/compiler-quartic-holonomy-20260919/PROOF.txt}\par{\footnotesize SHA-256 (Anfang): \texttt{5b5664c9704e9b60}.}\par
\subsection*{S13 — Korrelierte Ereignisse und Petersen-Clock}\label{s:13}
{\small\path{UR.COMPILER.CORRELATED_EVENT_CLOCK.13}\par}
Vertrag: \path{TFPT/experiments/theory-contracts/compiler-correlated-event-clock-20260919/contract_index.json}\par{\footnotesize SHA-256 (Anfang): \texttt{dcae87ea82c76b4f}.}\par
Herleitung: \path{TFPT/experiments/theory-contracts/compiler-correlated-event-clock-20260919/PROOF.txt}\par{\footnotesize SHA-256 (Anfang): \texttt{bf1f85c99f5f89a0}.}\par
\subsection*{S14 — Phasenbewusste Rahmengruppe}\label{s:14}
{\small\path{UR.COMPILER.FRAME_GROUP.14}\par}
Vertrag: \path{TFPT/experiments/theory-contracts/compiler-frame-group-20260919/contract_index.json}\par{\footnotesize SHA-256 (Anfang): \texttt{e70773cc99dd2a37}.}\par
Herleitung: \path{TFPT/experiments/theory-contracts/compiler-frame-group-20260919/PROOF.txt}\par{\footnotesize SHA-256 (Anfang): \texttt{60969a33006b3134}.}\par
\subsection*{S15 — Kontinuierlicher Ereignisursprung}\label{s:15}
{\small\path{UR.COMPILER.CONTINUOUS_EVENT_ORIGIN.15}\par}
Vertrag: \path{TFPT/experiments/theory-contracts/compiler-continuous-event-origin-20260919/contract_index.json}\par{\footnotesize SHA-256 (Anfang): \texttt{0f0f138705e6ed02}.}\par
Herleitung: \path{TFPT/experiments/theory-contracts/compiler-continuous-event-origin-20260919/PROOF.txt}\par{\footnotesize SHA-256 (Anfang): \texttt{545cb24fd942e67d}.}\par
\subsection*{S16 — Affiner Nachfahrenanschluss}\label{s:16}
{\small\path{UR.COMPILER.CURRENT_DESCENDANT.16}\par}
Vertrag: \path{TFPT/experiments/theory-contracts/compiler-current-descendant-20260919/contract_index.json}\par{\footnotesize SHA-256 (Anfang): \texttt{d27ef5b55d8716a6}.}\par
Herleitung: \path{TFPT/experiments/theory-contracts/compiler-current-descendant-20260919/PROOF.txt}\par{\footnotesize SHA-256 (Anfang): \texttt{ad45198c44b84365}.}\par
\subsection*{S17 — Quantisierte Quellenrückwirkung}\label{s:17}
{\small\path{UR.COMPILER.ROOT_SOURCE_BACKREACTION.17}\par}
Vertrag: \path{TFPT/experiments/theory-contracts/compiler-root-source-backreaction-20260919/contract_index.json}\par{\footnotesize SHA-256 (Anfang): \texttt{4df71c5f267eaf0f}.}\par
Herleitung: \path{TFPT/experiments/theory-contracts/compiler-root-source-backreaction-20260919/PROOF.txt}\par{\footnotesize SHA-256 (Anfang): \texttt{f61bf749daf63c10}.}\par
\subsection*{S18 — Phasentreue Lifts und Readouts}\label{s:18}
{\small\path{UR.COMPILER.ROOT_PHASE_LIFT.18}\par}
Vertrag: \path{TFPT/experiments/theory-contracts/compiler-root-phase-lift-20260919/contract_index.json}\par{\footnotesize SHA-256 (Anfang): \texttt{afb3140619c42796}.}\par
Herleitung: \path{TFPT/experiments/theory-contracts/compiler-root-phase-lift-20260919/PROOF.txt}\par{\footnotesize SHA-256 (Anfang): \texttt{d7465bd8ee3514d2}.}\par
\subsection*{S19 — Überlappende Quellen und starke Phasenkopplung}\label{s:19}
{\small\path{UR.COMPILER.OVERLAP_PHASE_LOCK.19}\par}
Vertrag: \path{TFPT/experiments/theory-contracts/compiler-overlap-phase-lock-20260919/contract_index.json}\par{\footnotesize SHA-256 (Anfang): \texttt{c43c5ee8545fd73a}.}\par
Herleitung: \path{TFPT/experiments/theory-contracts/compiler-overlap-phase-lock-20260919/PROOF.txt}\par{\footnotesize SHA-256 (Anfang): \texttt{e24e640c3c341e2f}.}\par
\subsection*{S20 — Ketten, räumliche Antwort und Stromvervollständigung}\label{s:20}
{\small\path{UR.COMPILER.SPATIAL_RESPONSE.20}\par}
Vertrag: \path{TFPT/experiments/theory-contracts/compiler-spatial-response-20260919/contract_index.json}\par{\footnotesize SHA-256 (Anfang): \texttt{a4913e3504199393}.}\par
Herleitung: \path{TFPT/experiments/theory-contracts/compiler-spatial-response-20260919/PROOF.txt}\par{\footnotesize SHA-256 (Anfang): \texttt{b200ef9a3d082cf3}.}\par
\subsection*{S21 — Native Stromprodukte}\label{s:21}
{\small\path{UR.COMPILER.CURRENT_PRODUCT.21}\par}
Vertrag: \path{TFPT/experiments/theory-contracts/compiler-current-product-20260919/contract_index.json}\par{\footnotesize SHA-256 (Anfang): \texttt{1a2843f7bf358abb}.}\par
Herleitung: \path{TFPT/experiments/theory-contracts/compiler-current-product-20260919/PROOF.txt}\par{\footnotesize SHA-256 (Anfang): \texttt{940131acfceba193}.}\par
\subsection*{S22 — Clock-erzwungene E8-Stromschließung}\label{s:22}
{\small\path{UR.COMPILER.CLOCK_CURRENT_CLOSURE.22}\par}
Vertrag: \path{TFPT/experiments/theory-contracts/compiler-clock-current-closure-20260919/contract_index.json}\par{\footnotesize SHA-256 (Anfang): \texttt{f06dd6132ad73df2}.}\par
Herleitung: \path{TFPT/experiments/theory-contracts/compiler-clock-current-closure-20260919/PROOF.txt}\par{\footnotesize SHA-256 (Anfang): \texttt{786817abb2707ee3}.}\par
\subsection*{S23 — Audit der Paar- und Mischsektoren}\label{s:23}
{\small\path{UR.COMPILER.PAIR_SECTOR_AUDIT.23}\par}
Vertrag: \path{TFPT/experiments/theory-contracts/compiler-pair-sector-audit-20260919/contract_index.json}\par{\footnotesize SHA-256 (Anfang): \texttt{49845a7badcaac7c}.}\par
Herleitung: \path{TFPT/experiments/theory-contracts/compiler-pair-sector-audit-20260919/PROOF.txt}\par{\footnotesize SHA-256 (Anfang): \texttt{b63d5a7a97c331f0}.}\par
\subsection*{S24 — Bindungsantwort und damalige Vollraumannahmen}\label{s:24}
{\small\path{UR.COMPILER.BOND_RESPONSE.24}\par}
Vertrag: \path{TFPT/experiments/theory-contracts/compiler-bond-response-20260919/contract_index.json}\par{\footnotesize SHA-256 (Anfang): \texttt{93c749362eda3d17}.}\par
Herleitung: \path{TFPT/experiments/theory-contracts/compiler-bond-response-20260919/PROOF.txt}\par{\footnotesize SHA-256 (Anfang): \texttt{67e3d788c0bfe616}.}\par
\subsection*{S25 — Domänenwandkern}\label{s:25}
{\small\path{UR.COMPILER.DOMAIN_WALL_CORE.25}\par}
Vertrag: \path{TFPT/experiments/theory-contracts/compiler-domain-wall-core-20260919/contract_index.json}\par{\footnotesize SHA-256 (Anfang): \texttt{35276acd6840b8a3}.}\par
Herleitung: \path{TFPT/experiments/theory-contracts/compiler-domain-wall-core-20260919/PROOF.txt}\par{\footnotesize SHA-256 (Anfang): \texttt{fe9c66ef6821ff47}.}\par
\subsection*{S26 — Relaxierter Kern und Zwei-Quellen-Vollsatz}\label{s:26}
{\small\path{UR.COMPILER.RELAXED_WALL.26}\par}
Vertrag: \path{TFPT/experiments/theory-contracts/compiler-relaxed-wall-20260919/contract_index.json}\par{\footnotesize SHA-256 (Anfang): \texttt{479ba50d2289f2ca}.}\par
Herleitung: \path{TFPT/experiments/theory-contracts/compiler-relaxed-wall-20260919/PROOF.txt}\par{\footnotesize SHA-256 (Anfang): \texttt{b4930e179bc5658c}.}\par
\subsection*{S27 — Relationale Inzidenz}\label{s:27}
{\small\path{UR.COMPILER.RELATIONAL_INCIDENCE.27}\par}
Vertrag: \path{TFPT/experiments/theory-contracts/compiler-relational-incidence-20260920/contract_index.json}\par{\footnotesize SHA-256 (Anfang): \texttt{40c82156b0826b7f}.}\par
Herleitung: \path{TFPT/experiments/theory-contracts/compiler-relational-incidence-20260920/PROOF.txt}\par{\footnotesize SHA-256 (Anfang): \texttt{d86af76b009c6126}.}\par
\subsection*{S28 — Vier Follow-ups und Vier-Quellen-Vollsatz}\label{s:28}
{\small\path{UR.COMPILER.FOUR_FOLLOWUPS.28}\par}
Vertrag: \path{TFPT/experiments/theory-contracts/compiler-four-followups-20260920/contract_index.json}\par{\footnotesize SHA-256 (Anfang): \texttt{f4937cd581079dd4}.}\par
Herleitung: \path{TFPT/experiments/theory-contracts/compiler-four-followups-20260920/PROOF.txt}\par{\footnotesize SHA-256 (Anfang): \texttt{ee0775ff2dbd37ac}.}\par
Herleitung: \path{TFPT/experiments/theory-contracts/compiler-four-followups-20260920/origin/PROOF.txt}\par{\footnotesize SHA-256 (Anfang): \texttt{82bf4dce1ff3a475}.}\par
Herleitung: \path{TFPT/experiments/theory-contracts/compiler-four-followups-20260920/vacuum/PROOF.txt}\par{\footnotesize SHA-256 (Anfang): \texttt{8ba73582b414db25}.}\par
Herleitung: \path{TFPT/experiments/theory-contracts/compiler-four-followups-20260920/large_chain/PROOF.txt}\par{\footnotesize SHA-256 (Anfang): \texttt{fde2c9f5d566a185}.}\par
\subsection*{S29 — Kritische Brücke, Rahmen und Fermi-Patch}\label{s:29}
{\small\path{UR.COMPILER.CRITICAL_BRIDGE.29}\par}
Vertrag: \path{TFPT/experiments/theory-contracts/compiler-critical-bridge-20260920/contract_index.json}\par{\footnotesize SHA-256 (Anfang): \texttt{ae730a091abc621a}.}\par
Herleitung: \path{TFPT/experiments/theory-contracts/compiler-critical-bridge-20260920/PROOF.txt}\par{\footnotesize SHA-256 (Anfang): \texttt{cede7d0ba0b115e1}.}\par
Herleitung: \path{TFPT/experiments/theory-contracts/compiler-critical-bridge-20260920/energy/PROOF.txt}\par{\footnotesize SHA-256 (Anfang): \texttt{9a27358119b20cce}.}\par
Herleitung: \path{TFPT/experiments/theory-contracts/compiler-critical-bridge-20260920/symmetry/PROOF.txt}\par{\footnotesize SHA-256 (Anfang): \texttt{87990f749b795ee6}.}\par
Herleitung: \path{TFPT/experiments/theory-contracts/compiler-critical-bridge-20260920/frames/PROOF.txt}\par{\footnotesize SHA-256 (Anfang): \texttt{1b7f093501662183}.}\par
Herleitung: \path{TFPT/experiments/theory-contracts/compiler-critical-bridge-20260920/origin/PROOF.txt}\par{\footnotesize SHA-256 (Anfang): \texttt{f9a1af5f48253a39}.}\par
\section{Spätere Quellen-, Rand-, Paar- und Registerarbeiten}
Die folgenden Registerpositionen ergänzen den frühen Sitzungsstrang. Alle aufgeführten Hauptverträge behalten das Gesamtverdict \texttt{PARTIAL}. Die konkreten exakten Sätze, bedingten Konstruktionen und Ausschlüsse werden in den zugehörigen Kapiteln getrennt dargestellt.
\subsection*{S30 — Geladene Sektoren und Projektionsgedächtnis}\label{s:30}
{\small\path{UR.CHARGED_SECTOR_MEMORY.01}\par}
Vertrag: \path{TFPT/experiments/theory-contracts/charged-sector-memory-gate-20260920/contract_index.json}\par
{\footnotesize SHA-256 (Anfang): \texttt{1ded5ac47e9d8d86}.}\par
Beleg: \path{TFPT/experiments/theory-contracts/charged-sector-memory-gate-20260920/PROOF.txt}\par
{\footnotesize SHA-256 (Anfang): \texttt{a69c18e858f001df}.}\par
\subsection*{S31 — Quellenspektrum, RG und gemeinsame Clockzeit}\label{s:31}
{\small\path{UR.SOURCE.RG_CLOCK_BRIDGE.01}\par}
Vertrag: \path{TFPT/experiments/theory-contracts/source-rg-clock-bridge-20260920/contract_index.json}\par
{\footnotesize SHA-256 (Anfang): \texttt{348d6b2834185ffb}.}\par
Beleg: \path{TFPT/experiments/theory-contracts/source-rg-clock-bridge-20260920/PROOF.txt}\par
{\footnotesize SHA-256 (Anfang): \texttt{58d870d69596b5bd}.}\par
\subsection*{S32 — Geladene Stromantwort und Viertelholonomie}\label{s:32}
{\small\path{UR.SOURCE.CHARGED_TIME.01}\par}
Vertrag: \path{TFPT/experiments/theory-contracts/charged-source-time-audit-20260920/contract_index.json}\par
{\footnotesize SHA-256 (Anfang): \texttt{6afc4925e2af026d}.}\par
Beleg: \path{TFPT/experiments/theory-contracts/charged-source-time-audit-20260920/PROOF.txt}\par
{\footnotesize SHA-256 (Anfang): \texttt{97ddf9206a860c00}.}\par
Beleg: \path{TFPT/experiments/theory-contracts/charged-source-time-audit-20260920/CAR_PROOF.txt}\par
{\footnotesize SHA-256 (Anfang): \texttt{d2c7a80cf0efce43}.}\par
Beleg: \path{TFPT/experiments/theory-contracts/charged-source-time-audit-20260920/HOLONOMY_REVIEW.txt}\par
{\footnotesize SHA-256 (Anfang): \texttt{12217b8c7e906e11}.}\par
\subsection*{S33 — Lokaler wachsender Quellengrenzwert}\label{s:33}
{\small\path{UR.SOURCE.LOCAL_LIMIT.01}\par}
Vertrag: \path{TFPT/experiments/theory-contracts/source-local-line-20260920/contract_index.json}\par
{\footnotesize SHA-256 (Anfang): \texttt{981ceb033cac3fc4}.}\par
Beleg: \path{TFPT/experiments/theory-contracts/source-local-line-20260920/SOURCE_LINE_PROOF.txt}\par
{\footnotesize SHA-256 (Anfang): \texttt{d527dd5bcb1769b6}.}\par
Beleg: \path{TFPT/experiments/theory-contracts/source-local-line-20260920/SOURCE_LINE_REVIEW.txt}\par
{\footnotesize SHA-256 (Anfang): \texttt{4812be09012e3af7}.}\par
Beleg: \path{TFPT/experiments/theory-contracts/source-local-line-20260920/LOCAL_LIMIT_REVIEW.txt}\par
{\footnotesize SHA-256 (Anfang): \texttt{0713338a5845f8ad}.}\par
Beleg: \path{TFPT/experiments/theory-contracts/source-local-line-20260920/POLARIZATION_GATE.txt}\par
{\footnotesize SHA-256 (Anfang): \texttt{920fb2434d404b71}.}\par
\subsection*{S34 — Rohträger und Ursprung des 3+2-Schnitts}\label{s:34}
{\small\path{UR.RAW_CARRIER_ORIGIN.01}\par}
Vertrag: \path{TFPT/experiments/theory-contracts/raw-carrier-origin-gate-20260920/contract_index.json}\par
{\footnotesize SHA-256 (Anfang): \texttt{9d9c50d4d571998f}.}\par
Beleg: \path{TFPT/experiments/theory-contracts/raw-carrier-origin-gate-20260920/PROOF.txt}\par
{\footnotesize SHA-256 (Anfang): \texttt{352dea875d7e3c89}.}\par
\subsection*{S35 — Gefahrener Flux und Carry im vollständigen See}\label{s:35}
{\small\path{UR.SOURCE.DRIVEN_CARRY.01}\par}
Vertrag: \path{TFPT/experiments/theory-contracts/source-driven-carry-20260920/contract_index.json}\par
{\footnotesize SHA-256 (Anfang): \texttt{cbbe468b9271ba7e}.}\par
Beleg: \path{TFPT/experiments/theory-contracts/source-driven-carry-20260920/PROOF.md}\par
{\footnotesize SHA-256 (Anfang): \texttt{f5aac7f0ac5d65ae}.}\par
Beleg: \path{TFPT/experiments/theory-contracts/source-driven-carry-20260920/FLUX_PROOF_REVIEW.txt}\par
{\footnotesize SHA-256 (Anfang): \texttt{3ac5f1f4a80852f5}.}\par
\subsection*{S36 — Quellenprobe und algebraischer Carry}\label{s:36}
{\small\path{UR.SOURCE.PROBE_CARRY.01}\par}
Vertrag: \path{TFPT/experiments/theory-contracts/source-probe-carry-20260920/contract_index.json}\par
{\footnotesize SHA-256 (Anfang): \texttt{446323cc83e29594}.}\par
Beleg: \path{TFPT/experiments/theory-contracts/source-probe-carry-20260920/PROOF.txt}\par
{\footnotesize SHA-256 (Anfang): \texttt{df420a7d9aef268e}.}\par
\subsection*{S37 — Randselektion und vollständiges Quellfunktional}\label{s:37}
{\small\path{UR.SOURCE.BOUNDARY_SELECTION.01}\par}
Vertrag: \path{TFPT/experiments/theory-contracts/source-boundary-selection-20260920/contract_index.json}\par
{\footnotesize SHA-256 (Anfang): \texttt{ae8c6138331415e9}.}\par
Beleg: \path{TFPT/experiments/theory-contracts/source-boundary-selection-20260920/PROOF.md}\par
{\footnotesize SHA-256 (Anfang): \texttt{4081e9f7bc19da9c}.}\par
Beleg: \path{TFPT/experiments/theory-contracts/source-boundary-selection-20260920/LANE_FUNCTIONAL.md}\par
{\footnotesize SHA-256 (Anfang): \texttt{66a87451a0bfddba}.}\par
\subsection*{S38 — Globale Ladungsgruppe und Flavor-Rang}\label{s:38}
{\small\path{UR.SOURCE.FLAVOR_ORIGIN.01}\par}
Vertrag: \path{TFPT/experiments/theory-contracts/source-flavor-origin-20260920/contract_index.json}\par
{\footnotesize SHA-256 (Anfang): \texttt{c5893a43e40f4848}.}\par
Beleg: \path{TFPT/experiments/theory-contracts/source-flavor-origin-20260920/PROOF.txt}\par
{\footnotesize SHA-256 (Anfang): \texttt{328158deea8bab9a}.}\par
\subsection*{S39 — Glue, Parität und lokale Quellfelder}\label{s:39}
{\small\path{UR.SOURCE.GRADED_LOCALITY.01}\par}
Vertrag: \path{TFPT/experiments/theory-contracts/source-graded-locality-20260920/contract_index.json}\par
{\footnotesize SHA-256 (Anfang): \texttt{59d28c4f50af4237}.}\par
Beleg: \path{TFPT/experiments/theory-contracts/source-graded-locality-20260920/PROOF.txt}\par
{\footnotesize SHA-256 (Anfang): \texttt{669309ea7f240397}.}\par
\subsection*{S40 — Energiepfad und Auswahl relevanter Nulloperatoren}\label{s:40}
{\small\path{UR.SOURCE.DYNAMICS_SELECTION.01}\par}
Vertrag: \path{TFPT/experiments/theory-contracts/source-dynamics-selection-20260920/contract_index.json}\par
{\footnotesize SHA-256 (Anfang): \texttt{d01538455203a3c0}.}\par
Beleg: \path{TFPT/experiments/theory-contracts/source-dynamics-selection-20260920/PROOF.txt}\par
{\footnotesize SHA-256 (Anfang): \texttt{03c09468b07cc68a}.}\par
